% Generated draft: requires mathematical and source review.
\documentclass[11pt]{article}
\usepackage[margin=25mm]{geometry}
\usepackage{fontspec}
\setmainfont{Latin Modern Roman}
\usepackage{amsmath,amssymb,mathtools,unicode-math}
\setmathfont{Latin Modern Math}
\usepackage{enumitem,needspace,xurl,hyperref}
\setlength{\parindent}{0pt}
\setlength{\parskip}{5pt}
\setlength{\emergencystretch}{3em}
\newcommand{\N}{\mathbb{N}}
\newcommand{\Z}{\mathbb{Z}}
\newcommand{\Q}{\mathbb{Q}}
\newcommand{\R}{\mathbb{R}}
\newcommand{\C}{\mathbb{C}}
\newcommand{\F}{\mathbb{F}}
\newcommand{\A}{\mathbb{A}}
\newcommand{\PP}{\mathbb{P}}
\newcommand{\E}{\mathbb{E}}
\newcommand{\eps}{\varepsilon}
\newcommand{\dd}{\,\mathrm{d}}
\begin{document}
\section*{UnsolvedMath ID: 701647}
\textbf{Source problem code:} catalogue entry
\subsection{Definitions and mathematical statement}
A finite group \(H\) is called soluble if its derived series, defined by \(H^{(0)}=H\) and \(H^{(i+1)}=[H^{(i)},H^{(i)}]\), reaches the trivial group. For a finite group \(H\), let \(F(H)\) denote its Fitting subgroup, the largest nilpotent normal subgroup of \(H\). Define \(F_0(H)=1\) and, recursively, let \(F_{i+1}(H)\) be the inverse image of \(F(H/F_i(H))\) under the quotient map. For soluble \(H\), its Fitting height is
\[
h(H)=\min\{r\in\Z_{\geq 0}:F_r(H)=H\}.
\]
In particular, \(h(1)=0\).

For a finite group \(G\) and a positive integer \(t\), put
\[
\mathcal M_t(G)=\{M<G:M\text{ is a maximal proper subgroup of }G,\ M\text{ is soluble, and }h(M)\leq t\}.
\]
Here maximality is among all proper subgroups of \(G\). Two members \(M,N\) of \(\mathcal M_t(G)\) are identified when \(N=gMg^{-1}\) for some \(g\in G\). Write \(c_t(G)\) for the number of these conjugacy classes.

The question is whether
\[
\forall t\in\Z_{\geq1}\ \exists n(t)\in\Z_{\geq1}\ \forall\text{ finite groups }G,\qquad
c_t(G)\geq n(t)\ \Longrightarrow\ \bigl(G\text{ is soluble and }h(G)\leq t\bigr).
\]
The supplied record labels this question open; that status is not independently verified here.
\subsection{Short English statement}
For each positive integer \(t\), does a sufficiently large number of conjugacy classes of soluble maximal subgroups of Fitting height at most \(t\) force a finite group itself to be soluble of Fitting height at most \(t\)? The threshold may depend on \(t\), but must work for every finite group.
\subsection{Sources}
\begin{itemize}
\item Local UnsolvedMath registry: permanent ID \(701647\), exact source code \texttt{catalogue entry}, title \emph{Kourovka Notebook Problem 19.99}. The supplied record has status \texttt{open} and upstream revision \texttt{d371ee917ac9647fbb06cd6fc529deec6f3dda82}.
\item Supplied source attribution: \emph{The Kourovka Notebook (Unsolved Problems in Group Theory)}, original problem 19.99, 19th issue (2018); the background also identifies the 21st issue edition (2026), \texttt{arXiv:1401.0300v48}. The supplied collection description names E. I. Khukhro and V. D. Mazurov as editors.
\item Proposed by A. F. Vasil'ev and T. I. Vasil'eva; the supplied proposed year is 2018. These bibliographic details and the supplied open status have not been independently verified.
\end{itemize}
\end{document}
